\section{Research Extension: Nuisance-Marginal Arrival-Time Design}
\label{sec:nuisance-design}
This extension, implemented after the earlier studies, replaces the diagnostic
information proxy with a numerically integrated information calculation for an
explicit, approximate observation model. It is an opt-in research mode, not an
efficacy claim. The established planner remains the default.

\subsection{Prior Art and Scope}
Nuisance-aware Bayesian experimental design is established: Bartuska et al.
derive small-noise approximations for nuisance marginalization
\cite{bartuska2022nuisance}. Robust information design under prior ambiguity is
also established \cite{go2022robust}; ARGUS does not implement that criterion or
inherit its guarantees. In the directly adjacent acoustic domain, Raorane et al.
optimize sensor-cluster placement for acoustic-emission source localization in
isotropic plates and include physical validation \cite{raorane2024placement}.
Consequently, acoustic localization combined with information-based sensing
cannot be claimed as a new principle here. The narrower contribution is an
inspectable implementation of live posterior-region compression,
nuisance-marginal scalar design, numerical diagnostics and paired evaluation
within the sequential bistatic-interrogation workflow. Its legal novelty and
commercial superiority have not been established.

\subsection{Mass-Preserving Spatial Compression}
Let $p_i$ denote posterior mass at cell position $\vect{x}_i$, expressed in
physical metres. Up to five spatial representatives define nearest-centre
regions $R_k$ using physical Euclidean distance. Unlike renormalizing only the
representative cells, the regional weights preserve the entire posterior:
\begin{equation}
 w_k=\sum_{i\in R_k}p_i,\qquad
 \sum_k w_k=1,\qquad q_{i|k}=p_i/w_k.
\end{equation}
For source $\vect{s}$ and receiver $\vect{r}$, define
$L_i=\|\vect{x}_i-\vect{s}\|+\|\vect{x}_i-\vect{r}\|$.
The region-conditioned expectation and variance of this bistatic path are
computed using $q_{i|k}$. With nuisance posterior means $b$ for timing offset and
$v$ for velocity, the approximate regional arrival distribution is
\begin{align}
Y\mid K=k,e &\sim \mathcal{N}(\mu_k,\sigma_k^2),\\
\mu_k &=b+\E_k[L]/v,\\
\sigma_k^2 &=\operatorname{Var}_k[L]/v^2+\sigma_{\rm meas}^2+\sigma_b^2\nonumber\\
 &\quad+\E_k[L^2]\sigma_v^2/v^4+\sigma_{{\rm pose},k}^2.
\end{align}
The last term propagates independent source/receiver placement second moments
through first-order path gradients, scaled by the actual panel dimensions.
Coincident probe/cell geometry uses a conservative finite sensitivity rather
than an undefined derivative. A scalar pose bias without a known direction is
included through its squared magnitude. These are assumptions, not measured
pose-error distributions.

The measurement floor is
$\sigma_{\rm meas}^2=1.5\times10^{-8}+(3/f_s)^2/3$ in seconds squared,
combining the existing model floor with an approximate matched-filter sampling
neighbourhood. It is not a learned timing calibration or a waveform-dependent
precision bound. Moment matching and first-order nuisance propagation make
this a surrogate for the real observation likelihood.

\subsection{Information Calculation and Numerical Audit}
For Gaussian density $f_k$, the design utility is
\begin{equation}
 I(K;Y\mid e)=\sum_k w_k\int f_k(y)
 \log_2\!\frac{f_k(y)}{\sum_jw_j f_j(y)}\,dy.
\end{equation}
It targets the regional latent variable $K$, not exact cell location or the
full waveform. Deterministic Gauss--Hermite quadrature uses 24 and 48 nodes,
increasing to 96 when their difference exceeds 0.002 bits. Log-sum-exp
evaluation avoids direct mixture underflow. The implementation reports raw
information, the clamp to $[0,H(K)]$, node count and last-order difference.
That difference is a numerical diagnostic, not a certified error bound.
Independent adaptive integration tests check a nontrivial mixture; other tests
check identical-response and separated-response limits, unit invariance,
physical-coordinate partitioning and non-mutating planning.

For compatibility with existing multiobjective weights, the legacy-named API
score becomes $I(K;Y\mid e)T/\max\{\log_2 K,1\}$, where $T$ is model trust.
Unscaled information is exposed separately in bits. The remaining cost,
waveform, repetition, constraint and baseline-guard terms are retained; this
is not a pure mutual-information maximizer. A nuisance penalty compares the
nominal and marginalized information, with the signed change also retained.
Region-dependent variances can themselves convey information, so arbitrary
heteroscedastic nuisance changes do not imply monotonic information loss.

\subsection{Frozen Paired Synthetic Evaluation}
The protocol reserved seeds 41000--41019 before evaluation, distinct from
pilot seeds 31000--31005. Each method receives exactly eight acquisitions and
24 shared candidate opportunities per step, independent of truth and the
method's posterior. Candidate and acquisition random streams are controlled
separately. The production inference and calibration controller remain active;
the planning horizon is one. A second context changes acquisition velocity by
$+2\%$, attenuation by $+4\%$ and noise by $\times2$, without giving these
changes to the inference model. All 80 runs completed, with no execution or
measurement-rejection failures. Source fingerprints were unchanged during
evaluation. Table~\ref{tab:nuisance-design} reports physical MAP error.

\begin{table}[t]
\centering
\caption{Held-out fixed-candidate synthetic ablation. Errors in mm;
20 seeds per context and method, eight acquisitions per run.}
\label{tab:nuisance-design}
\setlength{\tabcolsep}{3pt}
\small
\begin{tabular}{@{}llrrrr@{}}
\toprule
Context & Method & Mean & Median & P90 & $>30$ mm\\
\midrule
Nominal & Established & 15.12 & 13.85 & 23.80 & 1/20\\
Nominal & Nuisance & 15.10 & 15.20 & 20.43 & 1/20\\
Stress & Established & 180.14 & 191.93 & 285.29 & 17/20\\
Stress & Nuisance & 146.56 & 126.36 & 275.54 & 17/20\\
\bottomrule
\end{tabular}
\end{table}

Define paired advantage as established error minus nuisance-aware error.
The nominal mean advantage is 0.014 mm, with 95\% paired-bootstrap interval
$[-2.24,2.94]$ mm: three wins, seven losses and ten ties. Under stress it is
33.58 mm, interval $[-30.62,95.83]$ mm: ten wins and ten losses. Intervals use
10,000 paired-seed resamples. Both include zero; reliable superiority is not
established. Both methods exceed 30 mm in 85\% of stress cases. The new method
reaches 426.96 mm error on seed 41011, worse than the baseline by 220.46 mm.
No run crosses the predefined 0.70 false-confidence threshold, but stress
abstention is only 1/20 versus 0/20. Low thresholded false confidence does not
mean failures are recognized.

The supported outcome is an auditable observation-level design calculation,
not proven localization improvement or sim-to-real transfer. The method remains
explicitly experimental. The complete artifact preserves per-step actions,
errors, nuisance diagnostics, configuration and code hashes in
\path{research_results/nuisance_design_validation.json}. A compact in-app
report is generated from that artifact and hash-checked, rather than manually
transcribing favourable results. Subsequent tuning requires development data
and a new reserved confirmatory test set.

\subsection{Scientific Interface}
The redesigned interface represents the live posterior as a three-dimensional
relief whose height is probability relative to the current peak, not physical
defect depth. R90 contours are computed in millimetres from panel dimensions
and discrete posterior cell centres. They are not field-calibrated coverage.
Source/receiver markers reflect the recommended experiment; decorative moving
excitation rings are explicitly illustrative. Global pause, reduced-motion
support and off-screen animation suspension keep motion subordinate to the
scientific readout. The interface separates the old unitless utility, the new
model-information bits, posterior containment and actual benchmark error.
